\chapter*{Preliminaries}
\addcontentsline{toc}{chapter}{\protect\numberline{}{Preliminaries}}
%\section*{Preliminaries}
%\addcontentsline{toc}{section}{\protect\numberline{}{Preliminaries}}

%\paragraph{Alignment as shortest path}
%
%Alignment is equivalent to shortest path.

\paragraph{\dijkstra and \A}
\dijkstra's algorithm \citep{dijkstra1959note} finds a shortest path from $v_s$
to~$v_t$ by \emph{expanding} vertices in order of increasing distance $\g(u)$
from the start. The \A algorithm~\citep{hart1968formal,pearl1984heuristics},
instead, directs the search towards a target by expanding vertices in order of
increasing ${f(u) := g(u) + h(u)}$, where $h(u)$ is a heuristic function that
estimates the distance $\h(u)$ to the end and $g(u)$ is the shortest length of a
path from $v_s$ to $u$ found so far. A heuristic is \emph{admissible} if it is a
lower bound on the remaining distance, $h(u) \leq \h(u)$, which guarantees that
\A has found a shortest path as soon as it expands $v_t$. Heuristic $h_1$
\emph{dominates} another heuristic $h_2$ when $h_1(u) \ge h_2(u)$ for all vertices $u$.
A dominant heuristic will usually, but not always~\citep{holte2010common},
expand less vertices. Note that \dijkstra's algorithm is
equivalent to \A using a heuristic that is always $0$, and that both algorithms
require non-negative edge costs. Our variant of the \A algorithm is provided
in~\cref{GLOBALsec:astar}.

At a high level, \A maintains the set of all \emph{explored} states, initialized
with the set of sources $S$. Then, \A iteratively \emph{expands} the explored
state with lowest estimated cost by exploring all its neighbors, until it finds
a target. Here, the cost for node $u$ is estimated by the distance from source,
called $g(u)$, plus the estimate from the heuristic $h(u)$.

\paragraph{Heuristic functions for \A}

% paper: trie
%\subsection{Background: General \A algorithm} \label{TRIEsubsec:general-astar}
Given a weighted graph $G=(V,E)$ with $E \subseteq V \times V \times
\mathbb{R}_{\geq 0}$, the \A algorithm (abbreviated as \A) searches for the
shortest path from sources $S \subseteq V$ to targets $T \subseteq V$. It is an
extension of \dijkstra's algorithm that additionally leverages a \emph{heuristic
function} $h \colon V \to \mathbb{R}_{\geq 0}$ to decide which paths to explore
first.
%
If $h(u) \equiv 0$, \A is equivalent to \dijkstra's algorithm.

\paragraph{Heuristic function}
The heuristic function $h(u)$ estimates the
cost $h^*(u)$ of a shortest path in $G$ from $u$ to a target $t \in T$. Intuitively, a
good heuristic correlates well with the distance from $u$ to $t$.

To ensure that \A indeed finds the shortest path, $h$ should be
\emph{admissible}:

\begin{definition}[Admissible heuristic] A heuristic $h$ is \emph{admissible}
    if it provides a lower bound on the distance to the closest target: $\forall
    u. h(u) \leq h^*(u)$.
\end{definition}

While any admissible $h$ ensures that \A finds optimal
alignments~\cite{dechter_generalized_1985}, the specific choice of $h$
is critical for performance. In particular, decreasing the error $\delta(u) =
h^*(u)-h(u)$ can only improve the performance of
\A~\cite{dechter_generalized_1985}. Thus, a key contribution of ours is
a domain-specific heuristic $h$.

\paragraph{\A algorithm}
We aim to guarantee optimal alignment while optimizing the average runtime
to not reach its worst-case complexity. While \dijkstra is an algorithm that
explores graph nodes in the order of their distance from the start, \A is a
generalization of \dijkstra that also accounts for their distance from the
target. \A prioritizes the exploration of nodes that seem to be closer to the
target nodes. This way, \A can sometimes dramatically improve on the performance
of \dijkstra while remaining optimal.
%
In contrast, the heuristic function we introduce is more informative and
consistently outperforms \dijkstra.

\cref{alg:astar} shows a generic implementation of the \A algorithm,
roughly following~\cite{dechter_generalized_1985}.
We do not implement the reconstruction of the best alignment in order to simplify the presentation.
The procedure \mbox{\textsc{BacktrackPath}} traces the best alignment back to the $source$, based on remembered edges used to optimize $f$ for each alignment state.
%
\cref{alg:astar} also shows a simple implementation of \dijkstra in terms of \A.
We omit the implementation of \textsc{BacktrackPath} for simplicity.

\begin{algorithm}[t]
	\caption{\A~algorithm} \label{alg:astar}
	\begin{algorithmic}[1]
		\Function{\A}{$G\colon \text{Graph}$,
			$S\colon \text{Sources}$,
			$T\colon \text{Targets}$,
			$h\colon \text{Heuristic function}$}
		\State $g \gets \mli{Map}\colon (\text{Nodes} \to \mathbb{R}_{\geq 0})$
		\Comment Shortest paths lengths to explored nodes

		\State $f \gets \mli{Map}\colon (\text{Nodes} \to \mathbb{R}_{\geq 0})$
		\Comment $f(u)=g(u)+h(u)$ 

		\State $Q \gets \mli{MinPriorityQueue}(\mli{priority}=f)$ 
		\Comment Priorities according to $f$
		\ForAll{$s \in S$}
			\State $g[s] \gets 0.0,\, f[s] \gets 0.0$
			\State $Q.\mli{push}(s)$
			\Comment Initially, explore all $s \in S$
		\EndFor
		\While{$Q \neq \emptyset$}
			\State $\mli{curr} \gets Q.\mli{pop}()$
			\Comment Get state with minimal $f$ to be expanded
			\If{$\mli{curr} \in T$}
				\State \Return \Call{BacktrackPath}{$\mli{curr}$}
				\Comment Reconstruct a walk to $\mli{curr}$
			\EndIf
			\ForAll{$(\mli{curr},\mli{next},\mli{cost}) \in G.\mli{outgoingEdges}(\mli{curr})$}
			\State $g_\mli{next} \gets g[\mli{curr}] + \mli{cost}$
			\State $\hat{f}_\mli{next} \gets g_\mli{next} + h(\mli{next})$
				\Comment Candidate value for $f[\mli{next}]$
				\If{$\hat{f}_\mli{next} < f[\mli{next}{}]$}
					\State $g[\mli{next}] \gets g_\mli{next}$		
					\State $f[\mli{next}] \gets \hat{f}_\mli{next}$		
					\State $Q.\mli{push}(\mli{next})$
					\Comment Explore state $\mli{next}$
				\EndIf
		\EndFor
		\EndWhile
		\State \textbf{assert} $\mli{False}$
		\Comment Cannot happen if $T$ is reachable from $S$
		\EndFunction

		\Statex

		\Function{\dijkstra}{$G\colon \mli{Graph}$,
			$S\colon \mli{Sources}$,
			$T\colon \mli{Targets}$}
			\State $h(v) \gets 0.0$
			\Comment Constant-zero function $h$
			\State $\Call{\A}{G,S,T,h}$
		\EndFunction
	\end{algorithmic}
\end{algorithm}

% paper: seeds

%\subsection{\A~algorithm for finding a shortest path} \label{SEEDsec:astar}
%
The \A~algorithm is a shortest path algorithm that generalizes \dijkstra's
algorithm by directing the search towards the target.
Given a weighted graph $G=(V,E)$, the \A~algorithm finds a shortest path from
sources $S \subseteq V$ to targets $T \subseteq V$.
%
To prioritize paths that lead to a target, it relies on an admissible heuristic
function $h \colon V \to \mathbb{R}_{\geq 0}$, where $h(v)$ estimates the
remaining length of the shortest path from a given node $v \in V$ to a target
$t \in T$.

\paragraph{Algorithm}
% 
In a nutshell, the \A~algorithm maintains a set of \emph{explored} nodes,
initialized by all possible starting nodes $S$. It then iteratively
\emph{expands} the explored state $v$ with lowest estimated total cost $f(v)$ by
exploring all its neighbors. Here, $f(v) := g(v) + h(v)$, where $g(v)$ is the
distance from $s \in S$ to $v$, and $h(v)$ is the estimated distance from $v$ to
$t \in T$.
%
When the \A~algorithm expands a target node $t \in T$, it reconstructs the path
leading to $t$ and returns it.
%
\paragraph{Admissibility}
%
The \A~algorithm is guaranteed to find a shortest path if its heuristic $h$
provides a lower bound on the distance to the closest target, often referred to
as $h$ being \emph{admissible} or optimistic.

Further, the performance of the \A~algorithm relies critically on the choice of
$h$. Specifically, it is crucial to have low estimates for the optimal paths but
also to have high estimates for suboptimal paths.

\paragraph{Discussion}
%
To summarize, we use the \A~algorithm to find a shortest path from $\st{u}{0}$
to $\st{v}{|q|}$ in $\AG$. To guarantee optimality, its heuristic function
$h\st{v}{i}$ must provide a lower bound on the shortest distance from state
$\st{v}{i}$ to a terminal state of the form $\st{w}{\lvert q \rvert}$.
%
Equivalently, $h\st{v}{i}$ should lower bound the minimal cost of aligning
$q[i{:}]$ to $\RG$ starting from $v$, where $q[i{:}]$ denotes the suffix of $q$
starting at position $i$ ($0$-indexed).
%
The key challenge is thus finding a heuristic that is not only admissible but
also yields favorable performance.

% paper:global
% Shortest paths, \A for MSA and semi-global alignment (\astarix), gaps
\paragraph{Shortest paths and \A}
A pairwise alignment that minimizes edit distance corresponds to a shortest path
in the \emph{alignment graph}~\citep{vintsyuk1968speech,ukkonen1985algorithms}.
Assuming non-negative edit costs, a shortest path can be found using \dijkstra's
algorithm~\citep{ukkonen1985algorithms} (\cref{GLOBALfig1-dij}) or
\A~\citep{spouge1989speeding}. \A is an informed search algorithm which uses a
task-specific heuristic function to direct its search. Depending on the
heuristic function, a shortest path may be found significantly faster than by an
uninformed search such as \dijkstra's algorithm.

% paper: global
%\subsection{Seed heuristic}
%
%\dictum{Cut off your nose to spite your face.}
%\vskip 1em

\paragraph{Seed-and-extend}
Seed-and-extend is a commonly used paradigm for solving semi-global alignment
approximately~\citep{kucherov2019evolution}. Seeds are also used to define and
compute LCSk~\citep{benson2014longest}, a generalization of longest common
subsequence (LCS). In contrast, our \emph{\sh} speeds up finding an optimal
alignment by using seed matches to speed up the \A search. A limitation of the
existing \sh is the low tolerance to increasing error rates due to using only
long exact matches without accounting for their order.

Our seed heuristic uses seeds in a very different way than existing seeding
approaches: instead of searching for a good alignment around seed matches, it
punishes potential alignments by the lack of matches. This negation makes the
difference between finding a good alignment and proving that no other alignment
is better.

% paper-global
\paragraph{Paths, alignments, seeds and matches}
Any path from $\st ij$ to $\st{i'}{j'}$ in the alignment graph~$G$ represents a
\emph{pairwise alignment} (or just \emph{alignment}) of the substrings $A_{i
\dots i'}$ and $B_{j \dots j'}$. We denote with $d(u,v)$ the distance between
states $u$ and $v$. A shortest path $\pi^*$ corresponds to an optimal alignment,
thus $\cost(\pi^*) = d(v_s, v_t) = \ed(A, B)$. For a state $u$ we write $\g(u)
:= d(v_s, u)$ and $\h(u) := d(u, v_t)$ for the distance from the start to $u$
and from $u$ to the target $v_t$, respectively.

% paper-global
We outline algorithms for exact pairwise alignment and their fastest implementations for
biological sequences. Refer to~\citet{kucherov2019evolution} for approximate,
probabilistic, and non-edit distance algorithms and aligners.

\paragraph{Banding and bit-parallelization} When similar sequences are being
aligned, the whole DP table may not need to be computed. One such
output-sensitive algorithm is the \emph{banded} algorithm of
\citet{ukkonen1985algorithms} (\cref{GLOBALfig1-band}) which considers only states
near the diagonal within an exponentially increasing \emph{band}, and runs in
$\Oh(ns)$ time, where $s$ is the edit distance between the sequences. This
algorithm, combined with the \emph{bit-parallel optimization}
by~\citet{myers1999fast} is implemented by the \edlib
aligner~\citep{vsovsic2017edlib} that runs in $\Oh(ns/w)$ runtime, where $w$ is
the machine word size (nowadays 32 or 64).
% todo: cite Word RAM model: https://en.wikipedia.org/wiki/Word_RAM

\paragraph{Diagonal transition and WFA}
The $\Oh(ns)$ runtime complexity can be improved using the algorithm independently
discovered by \citet{ukkonen1985algorithms} and \citet{myers1986ano} that is
known as \emph{diagonal transition} \citep{navarro2001guided} (\cref{GLOBALfig1-wfa}).
It has an $\Oh(ns)$ runtime in the worst-case but only takes expected
$\Oh(n+s^2)$ time under the assumption that the input sequences are
random~\citep{myers1986ano}. This algorithm has been extended to linear and
affine costs in the \emph{wavefront alignment} (WFA)
algorithm~\citep{marco2021fast} in a way similar to~\citet{gotoh1982improved},
and has been improved to only require linear memory in
\wfa~\citep{marco2022optimal} by combining it with the \emph{divide and conquer}
approach of~\citet{hirschberg1975linear}, similar to \citet{myers1986ano}
algorithm for unit edit costs (\cref{GLOBALfig1-biwfa}).
% This new name, WFA, has since replaced the original \emph{diagonal transition}
% name and refers to both the unit cost and affine variants of the
% algorithm as well as the implementation.
Note that when each sequence letter has an error with a constant probability,
the total number of errors $s$ is proportional to $n$, so that even these
algorithms have a quadratic runtime.

\paragraph{Chains}
A state $u=\st ij\in V$ \emph{precedes} a state $v=\st {i'}{j'}\in V$, denoted
$u\preceq v$, when $i\leq i'$ and $j\leq j'$. Similarly, a match $m$ precedes a
match $m'$, denoted $m\preceq m'$, when the end of $m$ precedes the start of
$m'$. This makes the set of matches a partially ordered set.
A state $u$ precedes a match $m$ (denoted $u\preceq m$) when it precedes
the start of the match. A \emph{chain} of matches is a (possibly empty) sequence
of matches $m_1 \preceq \dots \preceq m_l$.

\paragraph{Contours}
To efficiently calculate maximal chains of matches, \emph{contours} are used.
Given a set of matches $\matches$, $\statescore(u)$ is the number of matches in
the longest chain $u\preceq m_0 \preceq \dots$, starting at $u$. The function
$\statescore\st ij$ is non-increasing in both $i$ and $j$. \emph{Contours} are
the boundaries between regions of states with $\statescore(u) = \ell$ and
$\statescore(u)<\ell$ (see~\cref{GLOBALfig:contours}). Note that contour $\ell$ is
completely determined by the set of matches $m\in \matches$ for which
${\statescore(\start(m)) = \ell}$~\citep{hirschberg1977algorithms}.

\citet{hunt1977fast} give an algorithm to efficiently compute $\statescore$ when
$\matches$ is the set of single-letter matches between $A$ and $B$, and
\citet{deorowicz2014efficient} give an algorithm when $\matches$ is the set of
$k$-mer exact matches.

% Output-sensitive overview
%\paragraph{Relaxations} Approximate
%algorithms~\citep{kucherov2019evolution}, an $\Oh(n \log n)$ algorithm that is
%exact with high probability assuming a limited error rate
%$e<3.485\%$~\citep[Proof of Lemma~25, p.~17]{ganesh2020near}. The
%asymptotically best know algorithm reaches
%$\Oh(n^2/\log^2n)$~\citep{masek1980faster} but is not practical.

In the following, we formalize the task of optimally aligning a read to a
reference graph in terms of finding a shortest path in an \emph{alignment
graph}.

\paragraph{Reference graph}
%
A reference graph $\RG=(\RGV,\RGE)$ encodes a collection of references to be
considered when aligning a read. Its directed edges $\RGE \subseteq \RGV \times
\RGV \times \Sigma$ are labeled by nucleotide letters from $\Sigma =
\{\texttt{A},\texttt{C},\texttt{G},\texttt{T}\}$, hence any walk
$\reference{\pi}$ in $\RG$ spells a string $\sigma(\reference{\pi}) \in
\Sigma^*$.

An alignment of a read $q \in \Sigma^*$ to a reference graph $\RG$ consists of
(i)~a walk $\reference{\pi}$ in $\RG$ and (ii)~a sequence of edits (matches,
substitutions, deletions, and insertions) that transform
$\sigma(\reference{\pi})$ to $q$. An alignment is \emph{optimal} if it minimizes
the sum of edit costs for a given real-valued cost model $\cedits = (\cmatch,
\csubst,\cdel, \cins)$.
%
We assume that edit costs are non-negative---a pre-requisite for the correctness
of \A. Further, we assume that $\cmatch \leq \csubst, \cins, \cdel$---a
prerequisite for the correctness of our heuristic.

We note that our approach naturally works for cyclic reference graphs.

\begin{figure}[t]
	\begin{alignat*}{20}
		(
			&\langle&& u &,& i   &\rangle&,
			&\langle&  v &,& i+1 &\rangle&,
			&&q[i],
			&&\cmatch
		&)&\in \AGE
		&& \quad \text{ if } (u,v,\ell) \in \RGE, \ell = q[i] & \qquad \text{(match)}\\
		%
		(
			&\langle&& u &,& i   &\rangle&,
			&\langle&  v &,& i+1 &\rangle&,
			&&q[i],
			&&\csubst
		&) &\in \AGE
		&& \quad \text{ if } (u,v,\ell) \in \RGE, \ell \neq q[i] & \qquad \text{(substitution)}\\
		%
		(
			&\langle&& u &,& i &\rangle&,
			&\langle&  v &,& i &\rangle&,
			&&\varepsilon,
			&&\cdel
		&) &\in \AGE
		&& \quad \text{ if } (u,v,\ell) \in \RGE & \qquad \text{(deletion)}\\
		%
		(
			&\langle&& u &,& i   &\rangle&,
			&\langle&  u &,& i+1 &\rangle&,
			&&q[i],
			&&\cins
		&) &\in \AGE
		&& \quad & \qquad \text{(insertion)},
	\end{alignat*}
	\caption[Formal definition of alignment graph]{Formal definition of
	alignment graph edges $\AGE \subseteq \AGV[q] \times \AGV[q] \times
	\Sigma_\varepsilon \times \mathbb{R}_{\geq 0}$. Here, $u,v \in \RGV$, $0
	\leq i < |q|$, $\ell \in \Sigma$, and $\varepsilon$ represents the empty
	string, indicating that letter $\ell$ was deleted.}
	\label{SEEDfig:graph-edges}
\end{figure}

\paragraph{Alignment graph}
%
In order to formalize optimal alignment as a shortest path finding problem, we
rely on an \emph{alignment graph} $\AG[q]=(\AGV[q],\AGE[q])$.
%
Its nodes $\AGV[q]$ are \emph{states} of the form $\langle v, i \rangle$, where
$v \in \RGV$ is a node in the reference graph and $i \in \{0, \dots, |q|\}$
corresponds to a position in the read $q$.
%
Its edges $\AGE[q]$ are selected such that any path $\alignment{}{\pi}$ in
$\AG[q]$ from $\langle u, 0 \rangle$ to $\langle v, i \rangle$ corresponds to an
alignment of the first $i$ letters of $q$ to $\RG$.
%
Further, the edges are weighted, which allows us to define an \emph{optimal
alignment} of a read $q \in \Sigma^*$ as a shortest path $\alignment{}{\pi}$ in
$\AG[q]$ from $\langle u, 0 \rangle$ to $\langle v, |q| \rangle$, for any $u, v
\in \RGV$.
%
\cref{SEEDfig:graph-edges} formally defines the edges $\AGE$.

% paper:global
\paragraph{Sequences}
The input sequences $A = \overline{a_0a_1\dots a_{n-1}}$ and $B =
\overline{b_0b_1 \dots b_{m-1}}$ are over an alphabet $\Sigma$ with $4$ letters.
We refer to substrings $\overline{a_i \dots a_{i'-1}}$ as $\substr Ai{i'}$, to
prefixes $\overline{a_0 \dots a_{i-1}}$ as $A_{<i}$, and to suffixes
$\overline{a_i \dots a_{n-1}}$ as $A_{\geq i}$. The \emph{edit distance}
$\ed(A,B)$ is the minimum number of insertions, deletions, and substitutions of
single letters needed to convert $A$ into $B$.

\paragraph{Alignment graph}
Let \emph{state} $\st{i}{j}$ denote the subtask of aligning the prefix $A_{<i}$
to the prefix $B_{<j}$. The \emph{alignment graph} (also called \emph{edit
graph}) $G(V,E)$ is a weighted directed graph with vertices $V = \{\st ij \vert
{0\leq i \leq n}, {0\leq j\leq m}\}$ corresponding to all states, and edges
connecting tasks to subtasks: edge ${\st ij \to \st{i{+}1}{j{+}1}}$ has cost $0$
if ${a_i = b_j}$ (match) and $1$ otherwise (substitution), and edges ${\st ij
\to \st{i{+}1}{j}}$ (deletion) and ${\st ij \to\st{i}{j{+}1}}$ (insertion) have
cost $1$. We denote the root state $\st 00$ by $v_s$ and the target state $\st
nm$ by $v_t$. For brevity we write $f(\st ij)$ as $f \st ij$. The alignment graph is
a natural representation of the alignment problem that provides a base for all
alignment algorithms.

We now describe the task of aligning a query to a reference graph. To this end,
we (i)~introduce the task of optimal alignment on a \emph{reference graph}, and
(iii)~introduce an alternative formulation in terms of an \emph{alignment
graph}, which is the basis for shortest path formulations of the optimal
alignment.
%
\cref{TRIEfig:graph-constructions} summarizes these different graph types.

\paragraph{Reference graph}
We encode the collection of references to which we want to align in a reference
graph, which captures genomic variation that a linear reference cannot
express~\cite{paten_genome_2017,garrison_variation_2018}.
%
We formalize a reference graph as a tuple $\RG=(\RGV,\RGE)$ of nodes $\RGV$ and
directed, labeled edges $\RGE \subseteq \RGV \times \RGV \times \Sigma$, where
the alphabet $\Sigma=\{\texttt{A},\texttt{C},\texttt{G},\texttt{T}\}$ represents
the four different nucleotides.
%
Note that in contrast to sequence graphs~\cite{rautiainen_aligning_2017}, we
label edges instead of nodes.

\paragraph{Path, spelling}
Any path $\pi=(e_1,\dots,e_k) \text { in } \RG$ induces a \emph{spelling}
$\sigma(\pi) \in \Sigma^*$ defined by $\sigma(e_1)\cdots\sigma(e_k)$, where
$\sigma(e_i)$ is the label of edge $e_i$ and $\Sigma^* := \bigcup_{k \in
\mathbb{N}} \Sigma^k$. We note that our approach naturally handles cyclic walks
and does not require cycle unrolling, a feature of
\bitparallel~\cite{rautiainen_bitparallel_2019} but missing from
\vg~\cite{garrison_variation_2018}, \pasgal~\cite{jain_accelerating_2019} and
\valigntool~\cite{kavya_sequence_2019}.

\paragraph{Alignment on reference graph}
An \emph{alignment} of \emph{query} $q \in \Sigma^*$ to a reference graph
$\RG=(\RGV,\RGE)$ consists of (i)~a path $\pi \text{ in } \RG$ and (ii)~a
sequence of edit operations (matches, substitutions, insertions, deletions)
transforming $\sigma(\pi)$ to~$q$.

\paragraph{Optimal alignment, edit Distance}
Each edit operation is associated with a real-valued cost ($\cmatch$, $\csubst$,
$\cins$, and $\cdel$, respectively).
An optimal alignment minimizes the total cost of the edit operations converting
$\sigma(\pi)$ to $q$. For optimal alignments, this total cost is equal to the
edit distance between $\sigma(\pi)$ and $q$, \ie, the cheapest sequence of edit
operations transforming $\sigma(\pi)$ into $q$.

We make the (standard) assumption that $0 \leq \cmatch \leq \csubst, \cins,
\cdel$, which will be a prerequisite for the correctness of our approach.

\begin{figure}[t]
	\centering
	\includegraphics[width=0.8\columnwidth]{edit_graph}
	\caption[Constructing the alignment graph]{Starting from the reference graph
	(left), we can construct the edit graph (middle) and the alignment graph
	$\AG$ for query $q=``\texttt{A}"$ (right). Edges are annotated with labels
	and/or costs, where sets of labels represent multiple edges, one for each
	letter in the set (indicated by $``\text{x}3"$ and $``\text{x}4"$).}
	\label{TRIEfig:graph-constructions}
\end{figure}

% TODO: remove
%\paragraph{Edit graph}
%Instead of representing alignments as pairs of (i)~paths in the reference graph and
%(ii)~sequences of edit operations on these paths, we introduce \emph{edit
%graphs} whose paths intrinsically capture both. This way, we can
%formally define an alignment more conveniently as a path in an edit graph.
%
%Formally, an \emph{edit graph} $\AG:=(\AGV,\AGE)$ has directed, labeled edges
%$\AGE \subseteq \AGV \times \AGV \times \Sigma_\epsilon \times \mathbb{R}_{\geq
%0}$ with associated costs that account for edits. Here, $\Sigma_\epsilon :=
%\Sigma \cup \{\epsilon\}$ extends the alphabet $\Sigma$ by $\epsilon$ to account
%for deleted characters (see \cref{TRIEfig:graph-constructions}).
%%
%The edit and reference graphs consist
%of the same vertices, \ie, $\AGV=\RGV$. However, $\AGE$ contains more edges
%than $\RGE$ to account for edits.
%%
%Concretely, for each edge $(u,v,\ell) \in \RGE$, $\AGE$ contains edges to
%account for (i)~matches, by an edge $(u,v,\ell,\cmatch)$, (ii)~substitutions, by
%edges $(u,v,\ell',\csubst)$ for each $\ell' \in \Sigma \backslash \ell$,
%(iii)~deletions, by an edge $(u,v,\epsilon,\cdel)$, and (iv)~insertions, by
%edges $(u,u,\ell',\cins)$ for each $\ell' \in \Sigma$.
%%
%The spelling $\sigma(\pi) \in \Sigma^*$ of a path $\pi \in \AG$ is defined
%analogously to reference graphs, except that deleted letters (represented by
%$\epsilon$) are ignored. The cost $\cost{\pi}$ of a path $\pi \in \AG$ is the
%sum of all its edge costs.
%
%\paragraph{Alignment on edit Graph}
%An \emph{alignment} of query $q$ to $\RG$ is a path $\pi \text{ in } \AG$
%spelling $q$, \ie, $q=\sigma(\pi)$. An \emph{optimal alignment} is an alignment
%of minimal cost.

\paragraph{Alignment graph}
To find an optimal alignment of $q$ using shortest path finding algorithms, we
must ensure that only paths spelling $q$ are considered. To this end, we
introduce an alternative but equivalent formulation of alignments in terms of an
\emph{alignment graph} $\AG=(\AGV,\AGE)$.

Here, each \emph{state} $\langle v,i \rangle \in \AGV$ consists of a vertex $v \in
\RGV$ and a query position $i \in \{0,\dots,|q|\}$ (equivalent
to~\cite{rautiainen_aligning_2017}). Traversing a state $\langle v,i \rangle \in
\AGV$ represents the alignment of the first $i$ query characters ending at node $v$.
%
In particular, query position $i=0$ indicates that we have not yet matched any
letters from the query.
%
We note that the alignment graph explicitly depends on the query $q$. In
particular, the example alignment graph $\AG[``\texttt{A}"]$ in
~\cref{TRIEfig:graph-constructions} lacks substitution edges from $\AG$, as their
labels ($\texttt{C}$, $\texttt{G}$, $\texttt{T}$) do not match the query
$q=``A"$.

%We construct the alignment graph $\AG$ to guarantee that any walk from a source
%$\langle u,0 \rangle$ to a state $\langle v,i \rangle$ corresponds to an
%alignment of the first $i$ letters of query $q$ to $\RG$.
%As a consequence, there is a one-to-one correspondence between alignments $\edit{\pi}$ of $q$ to
%$\AG$ and paths $\alignment{\pi} \in \AG$ from sources $S:=\RGV \times \{0\}$ to
%targets $T:=\RGV \times \{|q|\}$, with
%$\cost{\reference{\pi}}=\cost{\alignment{\pi}}$. To find the best alignment in
%$\AG$, only paths in $\AG$ (walks without repeating nodes) can be considered,
%since repeating a node in $\AG$ cannot lead to a lower cost ($\cdel \geq 0$) for
%the same state.

%The edges $\AGE \subseteq \AGV \times \AGV \times \Sigma_\epsilon \times
%\mathbb{R}_{\geq 0}$ are built based on the edges in $\AGE$, except that the
%former (i)~keep track of the position in the query $i$, and (ii)~only contain
%empty edges or edges
%whose label matches the next query letter:

\vspace{-1.2em}
{%
\small
\begin{alignat}{10}
	(u,v,\ell,w) &\in \AGE \implies (&\langle u, i \rangle, &\langle v, i+1
		&&\rangle,\ell,w) \in \AGE \quad \text{ for } 0 \leq i < |q| \text{ with }
		q[i]=\ell \label{TRIEeq:alignment-edges-nondeletions} \\
	(u,v,\epsilon,w) &\in\AGE \implies (&\langle u, i \rangle, &\langle v, i
		&&\rangle,\epsilon,w) \in \AGE \quad \text{ for } 0 \leq i < |q| \label{TRIEeq:alignment-edges-deletions}
\end{alignat}
}%

Here, assuming $0$-indexing, $q[i]$ is the next letter to be matched after
matching $i$ letters. Then, \cref{TRIEeq:alignment-edges-nondeletions} represents
matches, substitutions, and insertions (which advance the position in the query
by $1$), while \cref{TRIEeq:alignment-edges-deletions} represents deletions (which do
not advance the position in the query).